\section{Discussion}
\label{sec:discussion}

The revised results complicate any simple interpretation. It is not accurate to say that the model learned only output formatting and no mathematical relationship between input and output. Several diagnostics improved in ways that suggest the fine-tuned model captured broad numerical regularities associated with the inverse, including scale, sign structure, and reduced similarity to the input matrix itself. At the same time, those improvements remain insufficient for task success. Matrix inversion is an unforgiving functional target: small entrywise deviations can yield large residual errors, and coarse resemblance does not guarantee algebraic correctness. The results point to a gap between approximate structural alignment and exact functional computation.

\subsection{Why the Functional Gap Persists}

Several plausible, non-exclusive explanations may account for the remaining gap. These explanations are not fully separable without further ablation experiments.

\begin{enumerate}
    \item \textbf{Insufficient training signal.} 800 examples and 800 steps may be insufficient for a 3-billion-parameter model to generalize a precise function. Matrix inversion requires entrywise precision across all 4 or 9 output values; a small error in any entry can propagate into a large residual.

    \item \textbf{Representation mismatch.} LLMs process token sequences without intrinsic numeric semantics. Each floating-point value is tokenized as a sequence of character-level tokens. The model must implicitly treat these as numeric quantities and compose them correctly---a form of reasoning not supported by standard language modeling objectives.

    \item \textbf{Absence of intermediate computation.} Matrix inversion by Gaussian elimination proceeds through structured intermediate steps. Supervised fine-tuning on final ($A$, $A^{-1}$) pairs provides no signal about these steps, requiring the model to learn a direct input-to-output mapping that may be beyond what a small model can approximate from examples alone.

    \item \textbf{Selection effects in Level 2 diagnostics.} The structural improvements reported in Section~\ref{sec:experiments} are measured over the 60\% of fine-tuned outputs that are parseable. These may not be representative of the full fine-tuned distribution: outputs that terminated cleanly enough to be parsed may systematically differ from those that looped or became malformed. The apparent improvements could be overstated if loop-prone examples would have shown weaker numeric alignment.
\end{enumerate}

\subsection{Decoding Instability}

Format instability remained a primary bottleneck in this evaluation. All 100 fine-tuned outputs reached the 4{,}096-token inference ceiling, compared to typical baseline outputs of 52--55 tokens. The model lost clean length-controlled termination on every example after fine-tuning.

One possible explanation is that LoRA updates pushed the model toward generating more structured numeric content, and this effort leads to over-generation when the model cannot converge to a well-formed output---producing repeated blocks or fragmented shapes. This explanation is speculative; the precise cause is not established by the current evaluation.

The practical consequence is significant. A model that shows improved numeric structure among parseable outputs, but generates parseable output on only 60\% of examples, provides no net improvement in aggregate correctness (0 successes in both cases). Format instability must be addressed before the numeric improvements are meaningful in any applied sense.

\subsection{Interpretation}

A useful summary is to distinguish three levels of competence. First, a model may produce text that is syntactically well formed. Second, it may produce numbers that reflect some broad properties of the target distribution. Third, it may produce outputs that satisfy the mathematical function being asked. In this study, LoRA fine-tuning appears to improve the second level without reliably achieving the third, while also degrading robustness at the first level for 40\% of test examples. This result is informative: it demonstrates that the gap between structural resemblance and functional correctness is real and not closed by this training approach.

The results are compatible with a capacity-limitation interpretation---a sufficiently large model or more targeted training might succeed---but they do not establish it. The failure could equally reflect a more fundamental limitation of text-based arithmetic supervision.
